% Rabinovich 2018 Stavi presentation: exact operators and partition foundation.

\subsection{Temporal operators for proper gaps}
\label{sec:p09-stavi-gap-temporal-logic}

GPSS announces an arbitrary-order completeness theorem using
\(U',S'\), without giving their satisfaction clauses.
The distinct later presentation used here is Rabinovich's
\emph{A Proof of Stavi's Theorem}, in its revised published
2018 version. It defines two binary gap operators and
develops partition formulas for the full past/future
basis \cite[secs.~2--5]{Rabinovich2018Stavi}.
We write these operators as \(\mathsf U_{\mathrm{gap}}\)
and \(\mathsf S_{\mathrm{gap}}\). The notation
\(\mathsf U_s,\mathsf S_s\) continues to mean ordinary
strict until and since. No literal identification with
the undefined 1980 symbols is assumed.

\begin{hbdefinition}{def:p09-stavi-proper-gaps}{Proper gaps and the point/gap distinction}
In a nonempty strict linear order \(T\), a proper gap
is a partition \(T=L\mathbin{\dot\cup}R\) into
nonempty sets with every point of \(L\) below every
point of \(R\), where \(L\) has no greatest point
and \(R\) has no least point. Identify the gap
\(g\) with its downward-closed left side \(L\).
Write \(t<g\) when \(t\in L\), and \(g<t\)
when \(t\in R\). Intervals such as \((t,g)\)
and \((g,t)\) contain only actual points of \(T\).
For gaps \(g,h\), \(g<h\) means proper inclusion
of their left sides.

The source permits downward-closed proper cuts
whose supremum, if present, belongs to the cut.
Its proper gaps are exactly the two-sided gaps
just defined. A gap is not a temporal evaluation
point, and neither operand of a modality is
evaluated at a gap.
\end{hbdefinition}

\begin{hblemma}{thm:p09-stavi-gap-cuts}{proved=annex:p9.stavi.gap.cuts}{Gap boundaries and order classes}
The two definitions of proper gap agree.
A linear order has no proper gaps exactly
when it is Dedekind complete. If \(t<g\),
there is an actual point strictly between
\(t\) and \(g\); if \(g<t\), there is one
strictly between \(g\) and \(t\).
Distinct gaps have an actual point between
them. Order reversal exchanges their left
and right sides and preserves properness.
First and last points are not proper gaps.
\end{hblemma}

\begin{hbdefinition}{def:p09-stavi-definable-gap}{The gap determined by the first failure of a label}
A \(P\)-gap succeeding \(t\) is a proper gap
\(g>t\) such that \(P\) holds throughout
\((t,g)\) and
\[
        \forall r>g\ \exists u\,
                     (g<u<r\land\neg P(u)).
\]
Define the actual temporal formula
\[
\begin{split}
 \gamma^+(P)={}&
 (P\mathbin{\mathsf U_s}P)
 \land\neg(P\mathbin{\mathsf U_s}\neg P)
 \land F\neg P\\
 &{}\land\neg\bigl(P\mathbin{\mathsf U_s}
                        (P\land\partial^+\neg P)\bigr),
\end{split}
\]
where
\[
 \partial^+B(t)=
 \forall r\left(t<r\to
                    \exists u(t<u<r\land B(u))\right)
       =\neg((\neg B)\mathbin{\mathsf U_s}\top)(t).
\]
This is the vacuous right-limit convention.
The mirror formula \(\gamma^-(P)\) uses
strict since, strict past existence and
the corresponding left limit.
\end{hbdefinition}

\begin{hblemma}{thm:p09-stavi-definable-gap}{proved=annex:p9.stavi.definable.gap}{An explicit formula for the unique succeeding gap}
\(\gamma^+(P)\) holds at \(t\) exactly
when a \(P\)-gap succeeds \(t\), and that
gap is unique. The formula uses only
ordinary until and Boolean operations.
It tests \(P\) strictly after \(t\),
not at \(t\) itself. The mirror statement
holds for \(\gamma^-\).

For a first-order spelling put
\[
 \operatorname{Run}_P(t,u)=
          t<u\land\forall v(t<v<u\to P(v)).
\]
Then \(\Gamma_P^+(t)\), the first-order
translation of \(\gamma^+(P)\), is the
conjunction of
\[
\begin{gathered}
 \exists u\,(\operatorname{Run}_P(t,u)\land P(u)),\qquad
 \neg\exists u\,(\operatorname{Run}_P(t,u)\land\neg P(u)),\\
 \exists u\,(t<u\land\neg P(u)),\qquad
 \neg\exists u\,(\operatorname{Run}_P(t,u)\land
                              P(u)\land\partial^+\neg P(u)).
\end{gathered}
\]
All bound variables in its expanded constituents
are made fresh \cite[pp.~4--5]{Rabinovich2018Stavi}.
\end{hblemma}

\begin{hbdefinition}{def:p09-stavi-modalities}{The exact binary Stavi clauses}
The formula \(P\mathbin{\mathsf U_{\mathrm{gap}}}Q\)
holds at an actual point \(t\) exactly when
a \(P\)-gap \(g\) succeeds \(t\) and there
is an actual \(d>g\) such that \(Q\) holds
throughout \((g,d)\).
No condition is imposed on \(P(t)\), \(Q(t)\)
or \(Q(d)\).

The binary past dual
\(P\mathbin{\mathsf S_{\mathrm{gap}}}Q\)
holds at \(t\) exactly when a proper gap
\(g<t\) has \(P\) throughout \((g,t)\),
\(\neg P\) occurs arbitrarily close to
the left of \(g\), and some actual
\(d<g\) has \(Q\) throughout \((d,g)\).

The full language \(\mathcal T_{\mathrm{gap}}\)
is generated by proposition letters, Boolean
constants and operations, and the four
binary constructors
\(\mathsf U_s,\mathsf S_s,
\mathsf U_{\mathrm{gap}},\mathsf S_{\mathrm{gap}}\).
These definitions are the source's
\(\mathrm{Until}^{s},\mathrm{Since}^{s}\)
presentation, with different typography
to avoid confusing Stavi's superscript
with the handbook's strict-operator subscript.
\end{hbdefinition}

\begin{hbtheorem}{thm:p09-stavi-first-order-clauses}{proved=annex:p9.stavi.first.order.clauses}{First-order truth tables and all temporal constructors}
The future gap clause is exactly the
following first-order formula:
\[
\begin{split}
 \Gamma_P^+(t)\ \land\
 \exists x_1\bigl(t<x_1\land\neg P(x_1)\land
  \forall y\bigl(t<y<x_1\land\neg P(y)\to
                    \forall z(y<z<x_1\to Q(z))\bigr)\bigr).
\end{split}
\]
Writing \(\Gamma_P^-\) for the order
reverse of \(\Gamma_P^+\), the past
gap clause is explicitly
\[
\begin{split}
 \Gamma_P^-(t)\ \land\
 \exists x_1\bigl(x_1<t\land\neg P(x_1)\land
  \forall y\bigl(x_1<y<t\land\neg P(y)\to
                    \forall z(x_1<z<y\to Q(z))\bigr)\bigr).
\end{split}
\]
This reverses every order relation
without exchanging the two operands.
Together with the ordinary strict
since/until clauses, these give an
effective translation of every
\(\mathcal T_{\mathrm{gap}}\) formula
to monadic first-order logic, valid
on every nonempty linear order.

Both gap modalities are false on
Dedekind-complete orders. The future
one is false at a last point and
the past one at a first point.
Discreteness alone does not make
them false: an incomplete discrete
order can have a proper gap.
\end{hbtheorem}

\begin{hbdefinition}{def:p09-stavi-partitions}{Finite partitions and singleton flags}
For \(k\geq1\), labels \(D_1,\ldots,D_k\)
and \(O\subseteq\{1,\ldots,k\}\), write
\[
                  \operatorname{Part}(\langle D_1,\ldots,D_k\rangle,O)(I)
\]
when \(I\) is the union of \(k\) nonempty
convex intervals \(I_1<\cdots<I_k\), every
point of \(I_i\) satisfies \(D_i\), and
\(I_i\) is a singleton for each \(i\in O\).
An index outside \(O\) may also have a
singleton cell. Empty \(I\) cannot satisfy
a positive-length partition.

Initially the labels may be quantifier-free
one-variable formulas. In the canonical
construction they are supplied finite
\(\mathcal T_{\mathrm{gap}}\) formulas,
treated as names for their whole-model
truth sets and eventually inlined.
No interval endpoint is added as a free
parameter inside a label.

Basic partition formulas use \(x=y\),
\(x<y\), and these partitions on
\([x,y]\), \([x,+\infty)\),
\((-\infty,x]\), or the whole order.
Simple formulas are their finite
positive Boolean combinations.
Open or half-open bounded intervals
are expressed by padding the excluded
endpoints with true-labelled forced
singleton cells, with every original
flag shifted to its new position.
\end{hbdefinition}

\begin{hbdefinition}{def:p09-stavi-two-block}{A two-cell transition allowing a gap boundary}
Let \(P\mathbin{\mathsf U_*}Q\) abbreviate
\[
\begin{split}
 P\land\bigl[
    &(P\mathbin{\mathsf U_s}Q)
     \lor(Q\mathbin{\mathsf U_s}Q)\\
    &{}\lor
       P\mathbin{\mathsf U_s}
                   (P\land(Q\mathbin{\mathsf U_s}Q))
     \lor(P\mathbin{\mathsf U_{\mathrm{gap}}}Q)
                                                        \bigr].
\end{split}
\]
This is a finite formula in the declared
basis, not an additional primitive.
\end{hbdefinition}

\begin{hblemma}{thm:p09-stavi-two-block}{proved=annex:p9.stavi.two.block}{Every attained or missing two-cell boundary}
At \(t\), \(P\mathbin{\mathsf U_*}Q\)
holds exactly when some actual \(b>t\)
has a partition of \([t,b]\) into
two nonempty consecutive cells, labelled
\(P,Q\), with neither cell required
to be a singleton.
The four displayed alternatives cover
a first point of the second cell,
a last point of the first cell equal
to \(t\), a later last point of the
first cell, and a proper-gap boundary.
\end{hblemma}

\begin{hbdefinition}{def:p09-stavi-prefix-recursion}{The four-case finite partition recurrence}
For a labelled tuple and flags \(O\), put
\(F_k=D_k\), and for \(i=k,\ldots,2\) put
\[
 F_{i-1}=D_{i-1}\land
 \begin{cases}
    \bot\mathbin{\mathsf U_s}F_i,&i-1\in O,\ i\in O,\\
    D_i\mathbin{\mathsf U_s}F_i,&i-1\in O,\ i\notin O,\\
    D_{i-1}\mathbin{\mathsf U_s}F_i,&i-1\notin O,\ i\in O,\\
    D_{i-1}\mathbin{\mathsf U_*}F_i,&i-1\notin O,\ i\notin O.
 \end{cases}
\]
Every occurrence of \(\mathsf U_*\) is
expanded by its formula above.
For a future-ray translation replace
only the terminal seed by
\[
 F_k=
 \begin{cases}
   D_k\land G\bot,&k\in O,\\
   D_k\land GD_k,&k\notin O,
 \end{cases}
\]
retaining the original \(D_i\) in the
recurrence's interval operands.
\end{hbdefinition}

\begin{hbtheorem}{thm:p09-stavi-partition-prefix-rays}{proved=annex:p9.stavi.partition.prefix.rays}{Complete finite-prefix and ray translations}
With the finite seed, \(F_1(t)\) is
equivalent to the existence of \(b\geq t\)
such that the given partition holds
on \([t,b]\). More generally its suffix
formula \(F_i\) holds at every point
of the \(i\)-th cell of any realizing
partition, and expresses the corresponding
finite suffix starting at that point.

With the ray seed, \(F_1(t)\) is
equivalent to the partition on
\([t,+\infty)\), including a last
point and every singleton-flag case.
Reversing the order, tuple and flags,
and dualizing its labels,
gives the full past-ray translation.
These equivalences hold on every
linear order, with all labels evaluated
in the original whole model.
\end{hbtheorem}

\begin{hbtheorem}{thm:p09-stavi-positive-partitions}{proved=annex:p9.stavi.positive.partitions}{Finite refinement, anchoring and existential assembly}
Simple partition formulas are effectively
closed under finite conjunction,
disjunction and existential quantification.
The construction refines overlapping
partitions by their nonempty cell
intersections, enumerates weak order
types of named points, and splits
cells at each named point.

Each actual endpoint is represented
by its own forced singleton when
two distinct endpoints must be
retained. Existential elimination
then glues the two partitions at
the quantified singleton and
conjoins its labels. This preserves
the strict inequalities to surviving
endpoints. Every original singleton
flag, empty intersection case and
coincident named-point case is retained.
This proves the positive closure
needed before the source's negation
proposition.
\end{hbtheorem}

\begin{hblemma}{thm:p09-stavi-basic-partition-translations}{proved=annex:p9.stavi.basic.partition.translations}{Basic first-order definitions and the one-variable return}
Every basic partition formula has an
effective first-order presentation.
For a bounded closed interval, translate
its future-ray formula to first-order
logic and restrict its quantifiers
to that interval, while keeping the
input labels as unary names until
after restriction.

Every simple partition formula with
at most one free individual variable
has an effective
\(\mathcal T_{\mathrm{gap}}\) expression.
For a whole-order partition, choose
an actual \(z\) in its first cell,
combine its one-cell past-prefix
translation with its full future-ray
translation, and assert that this
combined formula holds somewhere:
\[
                               A\lor FA\lor PA.
\]
The past one-cell flag set is
\(O\cap\{1\}\). This restriction
retains a forced singleton first
cell without using an out-of-range
flag.
\end{hblemma}

\begin{hbdefinition}{def:p09-stavi-common-gap}{The same definable gap for two labels}
Define
\[
 \operatorname{UG}(D,E)=
     \gamma^+(D)\land\gamma^+(D\land E)
              \land\neg((D\land E)\mathbin{\mathsf U_{\mathrm{gap}}}D).
\]
\end{hbdefinition}

\begin{hblemma}{thm:p09-stavi-common-gap}{proved=annex:p9.stavi.common.gap}{The exact gap-capped interval test}
\(\operatorname{UG}(D,E)\) holds at \(t\)
exactly when a \(D\)-gap \(g\) succeeds
\(t\) and \(E\) holds throughout
\((t,g)\). The formula imposes no
condition on \(E(t)\).
The proof distinguishes equality of
the \(D\)-gap and \((D\land E)\)-gap
from a strictly earlier latter gap.
\end{hblemma}

\begin{hbdefinition}{def:p09-stavi-gap-partition}{The corrected open and closed gap-partition translations}
For \(k\geq1\), define
\(\operatorname{OG}_D(\langle E_1,\ldots,E_k\rangle,O)\).
If \(k\in O\), its value is false.
Otherwise form the tuple
\[
 \left\langle
   \top,\ D\land E_1,\ldots,D\land E_{k-1},
                D\land E_k\land\operatorname{UG}(D,E_k)
 \right\rangle
\]
with flags
\(\{1\}\cup\{i+1:i\in O\}\).
Apply the finite-prefix recurrence
to this tuple, and conjoin
\(\gamma^+(D)\) at its initial
point. The first cell is the
explicit initial singleton and
does not require \(D\) there.

The closed version
\(\operatorname{CG}_D(\langle E_1,\ldots,E_k\rangle,O)\)
is \(E_1\) conjoined with the
disjunction of these permitted
alternatives:
\[
\begin{array}{c|l}
 1\notin O&
  \operatorname{OG}_D(\langle E_1,\ldots,E_k\rangle,O),\\
 k\geq2&
  \operatorname{OG}_D(\langle E_2,\ldots,E_k\rangle,
                         \{i-1:i\in O,\ i\geq2\}).
\end{array}
\]
An absent alternative contributes
false. The second alternative
makes the original first cell
the current point alone.
\end{hbdefinition}

\begin{hbtheorem}{thm:p09-stavi-gap-partition}{proved=annex:p9.stavi.gap.partition}{Full partitions ending at a proper definable gap}
\(\operatorname{OG}_D\) holds at \(t\)
exactly when a \(D\)-gap \(g\)
succeeds \(t\) and the specified
partition holds on \((t,g)\).
\(\operatorname{CG}_D\) expresses
the corresponding condition on
\([t,g)\).

A forced singleton last cell
is impossible at a proper gap.
The current point in the closed
version must satisfy its own
label, but \(D(t)\) is not
required by a succeeding gap.
Order reversal gives both
right-open and right-closed
past versions with the same
explicit flag transformations.
\end{hbtheorem}

\begin{hbexample}{def:p09-stavi-gap-examples}{Dense and discrete incomplete examples}
On \(\mathbb Q\), let \(g\) be
the cut at \(\sqrt2\), take
\(P=\{q:q<\sqrt2\}\) and
\(Q=\{q:q>\sqrt2\}\).
At \(t=0\),
\(P\mathbin{\mathsf U_{\mathrm{gap}}}Q\)
holds. Every operand is evaluated
at a rational point; the missing
irrational is only a description
of the cut. Changing \(P(0)\)
alone does not change this value.

On the discrete order
\(\mathbb Z+\mathbb Z\), let
\(P\) be the first copy and
\(Q\) the second. The cut
between the copies is proper,
and the same gap formula holds
at every point of the first
copy. Thus local predecessor
and successor existence does
not eliminate proper gaps.
On a complete order both gap
modalities are false under
every interpretation.
\end{hbexample}

\begin{proofboundary}{def:p09-stavi-completeness-continuation}{The complete arbitrary-order construction}
For every monadic first-order formula \(\varphi(x)\),
\cref{thm:p09-stavi-completeness} constructs an equivalent formula
of \(\mathcal T_{\mathrm{gap}}\) over every nonempty linear order,
following the full target of the source
\cite[Theorems~2.3 and 4.4]{Rabinovich2018Stavi}.
The defined modalities, positive partitions and temporal returns here
provide the inputs to that construction.

The complete canonical negation of open partitions is
\cref{thm:p09-stavi-open-partition-complement}. It includes the
plain labelled-point-chain complement and the bounded prefix reduction,
resolves forced endpoint singleton cells before uniform-cell extension,
and retains all actual-boundary, ordered-gap and equal-gap cases.
\Cref{thm:p09-stavi-simple-negation} supplies Boolean negation closure;
the final theorem supplies every first-order induction step.
All gap formulas use the defined 2018 operators and retain the same
named cut whenever their conjuncts refer to a common missing point.
\end{proofboundary}
